\section{DPO-versus-SFT Results}
\label{sec:results}

For each DPO track, I compared the submitted checkpoint with the SFT checkpoint it started from, on the same 300 prompts, with the same generation settings and the same judge (Table~\ref{tab:dpo}). I report each comparison separately and never average across them. The two tracks use different prompts, rubrics, and preference data, so a single headline score would mix unlike targets.

\begin{table}[!t]
\centering
\caption{Pairwise judge results for each DPO model against SFT on identical prompts: wins (W), ties (T), losses (L), unknown verdicts (Unk.), and judge calls beyond one per batch of five (extra calls). Win rate excluding ties is W/(W+L) with a 95\% Wilson interval, win rate with ties counts each tie as half a win, and the sign test is exact and two-sided on W against L. The plain rubric drops the length guidance, and the repetition penalty applies to both models.}
\label{tab:dpo}
\begingroup\tablefont\setlength{\tabcolsep}{3pt}\renewcommand{\arraystretch}{1.0}
\begin{tabularx}{\textwidth}{>{\raggedright\arraybackslash}X>{\centering\arraybackslash}p{0.7cm}>{\centering\arraybackslash}p{0.7cm}>{\centering\arraybackslash}p{0.7cm}>{\centering\arraybackslash}p{0.7cm}>{\centering\arraybackslash}p{0.8cm}>{\centering\arraybackslash}p{1.0cm}>{\centering\arraybackslash}p{3.0cm}>{\centering\arraybackslash}p{1.4cm}>{\centering\arraybackslash}p{1.2cm}}
\toprule
Comparison & n & W & T & L & Unk. & Extra calls & Win rate, ties excluded [95\% CI] & Win rate, ties $\tfrac{1}{2}$ & Sign test $p$ \\
\midrule
\rowcolor{accent!8}\multicolumn{10}{@{}l}{\textbf{Own-track comparison: submitted DPO checkpoint against its SFT baseline}} \\
HH-RLHF DPO vs SFT, HH-eval & 300 & 137 & 70 & 93 & 0 & 8 & 0.596 [0.531, 0.657] & 0.573 & 0.004 \\
UltraFeedback DPO vs SFT, AlpacaEval & 300 & 94 & 60 & 146 & 0 & 1 & 0.392 [0.332, 0.455] & 0.413 & 0.001 \\
\rowcolor{accent!8}\multicolumn{10}{@{}l}{\textbf{Alignment tax: the other track's prompts}} \\
HH-RLHF DPO vs SFT, AlpacaEval & 300 & 102 & 93 & 105 & 0 & 2 & 0.493 [0.425, 0.560] & 0.495 & 0.889 \\
UltraFeedback DPO vs SFT, HH-eval & 300 & 96 & 68 & 134 & 2 & 20 & 0.417 [0.355, 0.482] & 0.436 & 0.015 \\
\rowcolor{accent!8}\multicolumn{10}{@{}l}{\textbf{Sensitivity of the own-track comparisons (HH-RLHF DPO on HH-eval, UltraFeedback DPO on AlpacaEval): one setting changed}} \\
HH-eval, positions swapped & 300 & 147 & 60 & 93 & 0 & 4 & 0.613 [0.550, 0.672] & 0.590 & <0.001 \\
HH-eval, plain rubric & 300 & 134 & 68 & 97 & 1 & 12 & 0.580 [0.516, 0.642] & 0.562 & 0.018 \\
HH-eval, repetition penalty 1.1 & 300 & 133 & 71 & 96 & 0 & 3 & 0.581 [0.516, 0.643] & 0.562 & 0.017 \\
AlpacaEval, positions swapped & 300 & 94 & 75 & 131 & 0 & 3 & 0.418 [0.355, 0.483] & 0.438 & 0.016 \\
AlpacaEval, plain rubric & 300 & 97 & 65 & 138 & 0 & 4 & 0.413 [0.352, 0.477] & 0.432 & 0.009 \\
AlpacaEval, repetition penalty 1.1 & 300 & 121 & 57 & 122 & 0 & 4 & 0.498 [0.436, 0.560] & 0.498 & 1.000 \\
\bottomrule
\end{tabularx}
\endgroup

\end{table}

\paragraph{Results.} On its own track, HH-RLHF DPO beat SFT on the HH-eval prompts, with 137 wins, 70 ties, and 93 losses. That is a win rate of 0.596 excluding ties (95\% interval 0.531 to 0.657) and 0.573 counting ties as half a win. UltraFeedback DPO lost to SFT on AlpacaEval, with 94 wins, 60 ties, and 146 losses, or 0.392 (0.332 to 0.455) and 0.413. Both gaps pass the sign test ($p=0.004$ and $0.001$). On the other track's prompts, HH-RLHF DPO matched SFT on AlpacaEval (102 wins, 93 ties, 105 losses; 0.493, $p=0.889$) and gave exactly SFT's answer on 30 of the 300 prompts, while UltraFeedback DPO lost on HH-eval (96, 68, 134; 0.417, $p=0.015$).

\paragraph{Judging procedure.} The judge is google/gemini-3.6-flash. It saw five prompts per request, with the two responses to each prompt in a seeded random A/B order (each model came first in exactly 150 of the 300 pairs) and no model names, and it answered A, B, or tie. The rubric matches the track: helpfulness and correctness for AlpacaEval, with a safety addendum for HH-eval. Both rubrics tell the judge not to prefer an answer for its length alone.

\paragraph{Failures and invalid verdicts.} I retried a failed request up to three times with a 5\,s pause, and then re-asked the pairs of a failed batch one at a time. A verdict that stayed invalid after that counts as unknown, and I excluded it from both win rates and logged it in a failure file. Only UltraFeedback DPO on the HH-eval prompts ended with unknown verdicts (2 of 300, after 3 failed batches and 20 extra calls), and a separate retry run did not resolve them. The plain-rubric check on HH-eval has one more. Pairs with identical outputs count as ties without a judge call, which covered 5 of the 300 UltraFeedback pairs and 2 of the 300 HH-RLHF pairs.

\paragraph{Sensitivity.} Swapping the A/B positions and removing the rubric's length guidance leave both conclusions in place: HH-RLHF scores 0.613 and 0.580, and UltraFeedback 0.418 and 0.413. A repetition penalty of 1.1, applied to both models, leaves HH-RLHF almost unchanged (133 wins, 71 ties, 96 losses; 0.581) but lifts UltraFeedback from 0.392 to 0.498 (121 wins, 57 ties, 122 losses), a tie ($p=1.0$). The UltraFeedback deficit therefore depends on the decoding setting, and the next section returns to why.

\paragraph{LoRA persona result.} The persona task has a different target, so I report it on its own (Table~\ref{tab:lorares}) and do not combine it with the DPO numbers. I judged each adapter against the base SmolLM2 model on persona style, usefulness, and safety. As a positive control, the judge compared the dataset's persona answers with its neutral answers on the 22 held-out prompts and picked the persona answer all 22 times, so the style rubric can see the persona. The selected adapter does not change the style the judge sees (0 wins and 40 ties on the required prompts, 2 wins and 20 ties on the held-out prompts), and its usefulness (16 wins, 3 ties, 21 losses) and safety (4, 32, 4) are not significantly different from the base model ($p=0.511$ and $1.0$). The strongest persona adapter, which used twice the learning rate and twice the epochs, wins on style (24 wins, 16 ties, 0 losses on the required prompts; 19, 3, 0 on the held-out prompts) but loses on usefulness (11, 2, 27; $p=0.014$) and safety (4, 18, 18; $p=0.004$). No LoRA verdict was unknown.

\begin{table}[!t]
\centering
\caption{LoRA persona results, each adapter against the base SmolLM2 model, one rubric per target. Table~\ref{tab:dpo} defines the columns.}
\label{tab:lorares}
\begingroup\tablefont\setlength{\tabcolsep}{3pt}\renewcommand{\arraystretch}{1.0}
\begin{tabularx}{\textwidth}{>{\raggedright\arraybackslash}X>{\centering\arraybackslash}p{0.7cm}>{\centering\arraybackslash}p{0.7cm}>{\centering\arraybackslash}p{0.7cm}>{\centering\arraybackslash}p{0.7cm}>{\centering\arraybackslash}p{3.0cm}>{\centering\arraybackslash}p{1.4cm}>{\centering\arraybackslash}p{1.2cm}}
\toprule
Judged target & n & W & T & L & Win rate, ties excluded [95\% CI] & Win rate, ties $\tfrac{1}{2}$ & Sign test $p$ \\
\midrule
\rowcolor{accent!8}\multicolumn{8}{@{}l}{\textbf{Positive control}} \\
Dataset persona answers vs neutral answers: style, 22 held-out prompts & 22 & 22 & 0 & 0 & 1.000 [0.851, 1.000] & 1.000 & <0.001 \\
\rowcolor{accent!8}\multicolumn{8}{@{}l}{\textbf{Selected adapter: rank 16, learning rate $5\times10^{-5}$, 3 epochs, NLL term (against the base model)}} \\
Persona style, 40 required prompts & 40 & 0 & 40 & 0 & - & 0.500 & - \\
Persona style, 22 held-out prompts & 22 & 2 & 20 & 0 & 1.000 [0.342, 1.000] & 0.545 & 0.500 \\
Usefulness, 40 required prompts & 40 & 16 & 3 & 21 & 0.432 [0.287, 0.591] & 0.438 & 0.511 \\
Safety, 40 required prompts & 40 & 4 & 32 & 4 & 0.500 [0.215, 0.785] & 0.500 & 1.000 \\
\rowcolor{accent!8}\multicolumn{8}{@{}l}{\textbf{Strongest persona adapter: rank 16, learning rate $1\times10^{-4}$, 6 epochs, NLL term (against the base model)}} \\
Persona style, 40 required prompts & 40 & 24 & 16 & 0 & 1.000 [0.862, 1.000] & 0.800 & <0.001 \\
Persona style, 22 held-out prompts & 22 & 19 & 3 & 0 & 1.000 [0.832, 1.000] & 0.932 & <0.001 \\
Usefulness, 40 required prompts & 40 & 11 & 2 & 27 & 0.289 [0.170, 0.448] & 0.300 & 0.014 \\
Safety, 40 required prompts & 40 & 4 & 18 & 18 & 0.182 [0.073, 0.385] & 0.325 & 0.004 \\
\bottomrule
\end{tabularx}
\endgroup

\end{table}
